\documentclass[11pt]{article}
\usepackage[T1]{fontenc}
\usepackage[margin=1in]{geometry}
\usepackage{amsmath,amssymb,amsthm}
\usepackage[hidelinks]{hyperref}
\newtheorem{theorem}{Theorem}
\title{An effective treewidth lower bound for Chromatic Number}
\author{Atlas\thanks{The proof in this note was found using ChatGPT.}}
\date{}
\begin{document}
\maketitle

Given a simple undirected graph \(G\) and a positive integer \(Q\),
Chromatic Number asks whether \(G\) admits a proper coloring with at most
\(Q\) colors. Let \(t\) be the treewidth of \(G\), let \(N\) be the binary
input length, and write \(s(t)=t\log_2(\max\{t,2\})\).

\hypertarget{tw-w1}{}\begin{theorem}\label{tw-w1}
Assume \(\mathrm{FPT}\ne\mathrm{W[1]}\). Chromatic Number has no
deterministic algorithm running in \(2^{h(t)}N^C\) time for a fixed
constant \(C\) and a finite nonnegative function \(h\) with the following
property: there is a computable function \(T\colon\mathbb{N}_{>0}\to
\mathbb{N}\) such that \(h(t)\le s(t)/j\) for every positive integer
\(j\) and every \(t\ge T(j)\).
\end{theorem}

\begin{proof}
Suppose such an algorithm exists. We give a fixed-parameter algorithm
for Clique parameterized by the requested clique size \(k\).
This problem is \(\mathrm{W[1]}\)-complete by graph complementation and
the corresponding theorem for Independent Set~\cite{df}.
Let the input graph have \(n\) explicitly represented vertices and binary
length \(M\). Handle \(k<2\) and \(k>n\) directly, and put
\(\ell=\lceil\log_2 n\rceil\).

Use \(k\) strings of \(\ell\) bits each to name the vertices of the
clique. In each string, forbid every unused name by a clause. For every
pair of strings, forbid each ordered pair of equal or nonadjacent
vertices by a clause that is false precisely when the two strings have
those names. The resulting formula has \(m=k\ell\) variables and
\(\mathcal O(kn+k^2n^2)\) clauses, each of length at most \(2\ell\).
It is satisfiable exactly when the input graph has a \(k\)-clique.

We next encode this formula by a coloring instance. For \(m<16\), test
all bit assignments directly. Otherwise put
\[
 b=\left\lfloor\tfrac12\log_2 m\right\rfloor,\qquad
 p=\lceil m/b\rceil,\qquad q=2^b .
\]
Split the variables into \(p\) groups of \(b\) bits, padding the last
group with unused bits. Make a clique on \(Q=q+p\) vertices to name
the colors: one assignment color \(A_a\) for each \(b\)-bit string
\(a\), and one witness color \(W_i\) for each group \(i\).
In a coloring with at most \(Q\) colors, all palette vertices have
distinct colors. In the construction below, we give a vertex an
allowed set of colors by joining it to every palette vertex outside
the set. This uses only ordinary edges.

For each group \(i\), introduce a vertex \(x_i\) allowed to use all
assignment colors. For each clause \(D\), introduce a selector
\(y_D\) allowed to use all witness colors. Choosing \(W_i\) means
that group \(i\) supplies a true literal of \(D\). To enforce this,
consider each group \(i\) and each assignment \(a\) that satisfies
no literal of \(D\) in that group. Add a private vertex adjacent to
\(x_i\) and \(y_D\), with allowed colors \(A_a\) and \(W_i\).
This also covers groups containing no literal of \(D\).

Each private vertex can be colored unless \(x_i\) has color \(A_a\)
and \(y_D\) has color \(W_i\). Private vertices are pairwise
nonadjacent. Consequently, given colors for the group vertices, a
clause selector and its private vertices can be colored exactly
when some group assignment makes a literal of that clause true.
The whole graph is thus \(Q\)-colorable exactly when the formula
is satisfiable.

Let \(S\) consist of the palette and all group vertices. A path
decomposition keeps \(S\) in every bag. Process the clauses in order:
use a bag containing \(S\) and the current selector, and add and
remove its private vertices one at a time. These bags contain both
ends of every edge, and the bags containing each vertex form an
interval. Thus the graph has treewidth \(t\le q+2p+1\).
Since \(q\le\sqrt m\) and \(b\ge(\log_2 m)/4\), this gives
\(t=\mathcal O(m/\log m)\). In particular, there is a fixed positive
integer \(A\) such that \(s(t)\le A m\) throughout the construction.
There are at most \(pq\) private vertices per clause. The graph,
including every edge that implements an allowed set, has polynomial
encoding length in \(M\) and can be constructed in polynomial time.

Finally, choose \(j=Ak\). If \(t\ge T(j)\), we have
\(h(t)\le A m/j=\ell\). To handle the other values of \(t\), let \(D_0\)
be a fixed nonnegative integer bounding all \(h(u)\) for \(u<T(1)\).
Such an integer exists since this set is finite. For all \(u\) we
then have \(h(u)\le D_0+s(u)\). Hence for \(t<T(j)\),
\(h(t)\le D_0+(T(j)+1)^2\), a computable function of \(k\):
the integer \(D_0\) is a constant of the assumed algorithm.
Since \(2^\ell\le 2n\), the resulting Clique algorithm runs in
\(f(k)M^{\mathcal O(1)}\) time for some computable \(f\).
This contradicts the hypothesis.
\end{proof}

\begin{thebibliography}{1}
\bibitem{df} R.~G. Downey and M.~R. Fellows.
Fixed-parameter tractability and completeness II: On completeness for W[1].
\emph{Theoretical Computer Science} 141(1--2), 109--131, 1995.
\href{https://doi.org/10.1016/0304-3975(94)00097-3}{DOI}.
\end{thebibliography}
\end{document}
